\begin{lemma} \label{lemma:accumulation_point_criteria_for_subsets_of_real_numbers}
    Let $\bbR$ be the \CrefAndHyperrefIfExist{definition:field_of_real_numbers}{field of real numbers} equipped with the \CrefAndHyperrefIfExist{definition:standard_absolute_value_on_real_numbers}{standard absolute value}. Let $A \subseteq \bbR$ and $p \in \bbR$.

    \begin{enumerate}
        \item $p$ is an \CrefAndHyperrefIfExist{definition:accumulation_point_in_an_extended_metric_space}{accumulation point} of $A$ if and only if for every $\varepsilon > 0$ there exists $a \in A$ such that $a \neq p$ and $|a - p| < \varepsilon$.
        \item $p$ is an accumulation point of $A$ if and only if there exists a \CrefAndHyperrefIfExist{definition:convergence_and_limit_of_sequence_in_extended_metric_space}{convergent sequence} $(a_n)_{n \in \bbN}$ in $A \setminus \{p\}$ such that $\lim_{n \to \infty} a_n = p$.
        \item Let $I$ be an \CrefAndHyperrefIfExist{definition:interval_in_the_real_line}{interval in the real line} that is \emph{not a singleton}\CrefIfExists{definition:singleton_set}. Suppose $p \in I \subseteq A$. Then $p$ is an accumulation point of $A$.
    \end{enumerate}
\end{lemma}
\begin{proof}
    \TODO{AI generated, verify}
    \begin{enumerate}
        \item This is exactly the $\varepsilon$-$\delta$ definition of an accumulation point in the extended metric space $\bbR$ as in \Cref{definition:accumulation_point_in_an_extended_metric_space}.

        \item Suppose $p$ is an accumulation point of $A$. For each $n \in \bbN$, the open ball $B(p, 1/(n+1))$ contains a point $a_n \in A \setminus \{p\}$ by part~(1). Then $|a_n - p| < 1/(n+1)$, so $(a_n)$ converges to $p$ in the sense of \Cref{definition:convergence_and_limit_of_sequence_in_extended_metric_space}.

        Conversely, suppose there exists a sequence $(a_n)_{n \in \bbN}$ in $A \setminus \{p\}$ converging to $p$. Let $\varepsilon > 0$. There exists $N \in \bbN$ such that $|a_n - p| < \varepsilon$ for all $n \geq N$. In particular, $a_N \in A \setminus \{p\}$ and $|a_N - p| < \varepsilon$, so $p$ is an accumulation point of $A$ by part~(1).

        \item Since $I$ is an interval that is not a singleton, $p$ is not isolated in $I$. Let $\varepsilon > 0$. We consider three cases.

        If $p$ is an interior point of $I$, there exists $\delta > 0$ such that $(p - \delta, p + \delta) \subseteq I$. Choose $a = p + \frac{\min(\delta, \varepsilon)}{2}$. Then $a \in I$, $a \neq p$, and $|a - p| < \varepsilon$.

        If $p$ is a left endpoint of $I$, then $I$ contains an interval $(p, b)$ for some $b > p$ because $I$ is not a singleton. Choose $a = p + \frac{\min(b - p, \varepsilon)}{2}$. Then $a \in I$, $a \neq p$, and $|a - p| < \varepsilon$.

        If $p$ is a right endpoint of $I$, the argument is symmetric: $I$ contains $(c, p)$ for some $c < p$, and we choose $a = p - \frac{\min(p - c, \varepsilon)}{2}$.

        In every case, condition~(1) is satisfied, so $p$ is an accumulation point of $A$.
    \end{enumerate}
\end{proof}